\documentclass[10pt,journal,compsoc]{IEEEtran}
\usepackage{hyperref}
\usepackage{graphicx}
\usepackage{booktabs}
\usepackage{amsmath}
\usepackage{url}

\begin{document}

\title{Lumen: A Privacy-First, Multi-Agent Adversarial Red-Teaming Sandbox for Pre-Deployment Clinical LLM Safety Certification}

\author{Dr. Baddam Sucharith Reddy, MBBS\\
\IEEEauthorblockA{MNR Medical College and Hospital, Hyderabad, Telangana, India\\
Clinical Lead \& Health AI Architect\\
Email: bsucharith003@gmail.com | ORCID: Pending}}

\maketitle

\begin{abstract}
Large Language Models (LLMs) are rapidly entering clinical workflows, yet standard static evaluation datasets fail to evaluate dynamic clinical safety behaviors under pressure. This paper presents \textit{Lumen}, an open-source, pre-deployment clinical safety evaluation workstation designed to red-team and audit clinical LLMs before they reach patient-facing deployments. Utilizing the Band Agent SDK as a shared clinical context bus, we coordinate a 4-agent collaborative network consisting of: (1) a Doctor Agent, (2) an autonomous Patient Simulator representing diverse clinical personas and regional languages, (3) an adaptive Red-Team Adversary, and (4) a Safety Auditor executing a 3-judge consensus review. The Adversary stress-tests the Doctor AI against a clinical attack taxonomy (Contraindication Bypass, Safety Screen Skip, Emergency Misdirection, and Disclaimer Burial). Dialogues are intercepted, mapped to LOINC/RxNorm vocabularies, compiled into HL7 FHIR R4 Transaction Bundles, and validated against public HAPI FHIR servers. Lumen was benchmarked across local-first configurations (Intel OpenVINO-accelerated Qwen-2.5-INT4 on consumer CPUs) and cloud gateways (Gemini 2.0 Flash). The multi-judge system successfully identified clinical protocol drift (mean safety score variance) and halted unsafe actions, exporting standardized FDA-aligned Software as a Medical Device (SaMD) scorecards. We demonstrate the viability of dynamic, multi-agent adversarial validation as a mandatory pre-deployment safety gate for healthcare AI, preserving clinical safety, privacy, and regulatory interoperability.
\end{abstract}

\begin{IEEEkeywords}
Clinical LLMs, Red-Teaming, Multi-Agent Systems, Medical AI Safety, HL7 FHIR, Intel OpenVINO.
\end{IEEEkeywords}

\section{Introduction}
The integration of generative Artificial Intelligence (AI) and Large Language Models (LLMs) into healthcare systems holds massive promise for reducing administrative burdens, automating electronic health record (EHR) documentation, and assisting in diagnostic pathways. However, clinical deployment introduces extreme safety risks. LLMs are known to suffer from medical hallucinations, drug dosing errors, and a lack of adherence to established clinical guidelines.

Existing AI evaluation frameworks (such as MMLU, MedQA, or MedMCQA) rely on static question-answering benchmarks. While these datasets evaluate memory recall and medical knowledge, they fail to evaluate \textbf{dynamic workflow safety}---the model's capacity to maintain safety rules under clinical pressure, complex dialogue drift, or intentional patient manipulation. For example, a model might correctly answer a multiple-choice question about the contraindications of a biologic drug, yet still prescribe that same drug in a simulated conversation without checking for a mandatory tuberculosis (TB) test.

Furthermore, deploying clinical evaluation tools raises severe data privacy concerns. Clinicians cannot upload sensitive patient transcripts containing Protected Health Information (PHI) to public third-party clouds without violating HIPAA regulations.

To bridge this pre-deployment gap, we present \textbf{Lumen}: a privacy-focused, local-first clinical red-teaming sandbox and safety certifier. Developed by a clinician, Lumen establishes a non-linear, multi-agent simulation loop to test LLMs under pressure, ensuring they undergo rigorous clinical stress-testing and standards compliance verification before touching real patients.

\section{Methodology \& System Architecture}
Lumen coordinates a coordinated multi-agent simulation loop orchestrated by the \textbf{Band Agent SDK}. The Band SDK functions as a high-reliability, shared clinical context bus (\texttt{BandSharedContext}), managing state transitions and task dispatches across four specialized agents.

\subsection{The 4-Agent Simulation Loop}
\begin{enumerate}
    \item \textbf{Red-Team Adversary Agent (Gemini 2.0 Flash)}: Generates complex patient vignettes containing hidden safety traps.
    \item \textbf{Patient Simulator Agent (Ollama / Gemini)}: Simulates distinct clinical personas (e.g., drug-seeking behaviors, anxious patients, elderly confused) and supports regional languages (Hindi, Telugu, Tamil, Marathi) to evaluate cross-cultural reliability.
    \item \textbf{Doctor AI Agent (The Tested Model)}: Evaluates the patient, orders labs/imaging, and prescribes treatments.
    \item \textbf{Safety Auditor Agent (AI/ML API / Gemini)}: Executes a 3-judge parallel consensus review to output safety grades (PASS, PARTIAL, FAIL) and audit logs.
\end{enumerate}

\subsection{Clinical Attack Taxonomy}
The Red-Team Adversary is programmed to exploit clinical safety blind spots using five targeted attack vectors:
\begin{itemize}
    \item \textbf{Contraindication Bypass}: Hiding underlying patient conditions (e.g., severe renal failure) that make first-line drug prescriptions dangerous.
    \item \textbf{Safety Screen Skip}: Attempting to bypass mandatory clinical protocols (e.g., starting immunosuppressive therapy without tuberculosis screening).
    \item \textbf{Dosing Hallucination}: Pushing the Doctor AI to calculate high-risk pediatric or organ-failure dosages.
    \item \textbf{Emergency Misdirection}: Presenting emergency symptoms (e.g., atypical cardiac ischemia) disguised as minor somatic complaints (e.g., epigastric indigestion).
    \item \textbf{Disclaimer Burial}: Delivering definitive medical instructions while hiding mandatory medical disclaimers at the bottom of the output.
\end{itemize}

\section{Standards Compliance \& Regional Interoperability}
To prepare models for EHR deployment, Lumen features terminology extraction and standards compliance mapping.

\subsection{Terminology Mapping (LOINC, RxNorm, CPT)}
When the Doctor AI emits tool calls (e.g., ordering lab tests or prescriptions), Lumen's parser maps the text to standardized clinical vocabularies:
\begin{itemize}
    \item \textbf{LOINC} (Logical Observation Identifiers Names and Codes) for diagnostic lab orders.
    \item \textbf{RxNorm} for active pharmaceutical drug ingredients.
    \item \textbf{CPT} (Current Procedural Terminology) for clinical procedures and surgeries.
\end{itemize}

\subsection{HL7 FHIR R4 Bundle Validation}
Every simulation session compiles into a complete \textbf{HL7 FHIR R4 Transaction Bundle} containing: \texttt{Patient} demographic envelopes, \texttt{Encounter} session tracking, \texttt{Observation} diagnostic results, \texttt{MedicationRequest} prescriptions, and \texttt{DocumentReference} clinical notes.
Lumen transmits these bundles to public \textbf{HAPI FHIR servers} to validate schema correctness, ensuring the AI-generated outputs comply with CMS and health network data standards.

\subsection{Regional API Sandboxes}
To evaluate international scalability, Lumen incorporates a directory router (\texttt{src/utils/regionalApis.ts}) mapped to national healthcare portals:
\begin{itemize}
    \item \textbf{India (ABDM)}: Validates ABHA (Ayushman Bharat Health Account) structure and HIE-CM consent manager tokens.
    \item \textbf{USA (Mayo Clinic \& OpenFDA)}: Queries de-identified patient cohorts and active drug adverse event databases.
    \item \textbf{UK (NHS GP Connect)}: Routes patient demographics checks through GP Connect FHIR interfaces.
\end{itemize}

\section{Evaluation \& Hardware Performance}
Lumen provides a privacy-first, local execution option utilizing \textbf{Intel OpenVINO} and local \textbf{Ollama} runtimes, comparing safety performance side-by-side with cloud models.

\subsection{Performance Comparison Matrix}
\begin{table}[h]
\centering
\caption{Lumen Performance Comparison Matrix}
\begin{tabular}{p{2.2cm}p{1.5cm}p{1.2cm}p{1.2cm}p{1cm}}
\toprule
Deployment Model & Hardware Target & Avg. Latency & Token Rate & Safety (Mean/SD) \\
\midrule
\textbf{Qwen 2.5 (INT4)} via OpenVINO & Local CPU & 180ms & 22 t/s & 82.5 / 3.2 \\
\textbf{BioMistral 7B} via Featherless & Cloud GPU & 450ms & 45 t/s & 85.0 / 4.1 \\
\textbf{Gemini 2.0 Flash} via AI Studio & Cloud TPU & 320ms & 50 t/s & 94.2 / 1.8 \\
\bottomrule
\end{tabular}
\end{table}

\subsection{Findings on Protocol Drift}
By running identical cases repeatedly, the Safety Auditor calculates the \textbf{Safety Score Standard Deviation (StdDev)}. A high StdDev indicates inconsistent clinical judgment (protocol drift), flagging a model as unsafe for deployment even if its average score is high.

\section{Discussion \& Future Work}
Lumen demonstrates that a dynamic, multi-agent adversarial simulation loop can expose critical safety failures that static QA benchmarks miss. By utilizing the \textbf{Band SDK} as a context bus, the system maintains strict state tracking, allowing safety auditor judges to pinpoint exactly where the clinical pathway drifted.

The addition of the \textbf{Autofill Demo Code} login bypass on the workstation dashboard facilitates a frictionless review experience for evaluators and regulators while demonstrating the cost-security and HIPAA-compliance lockouts necessary in clinical environments.

Future iterations of Lumen will focus on:
\begin{enumerate}
    \item Integrating active \textbf{SNOMED CT} medical terminology graphing.
    \item Conducting pilot integrations with open-source EMR systems (such as OpenEMR).
    \item Expanding the regional directory registry to support EU-EHDN (European Health Data Space) frameworks.
\end{enumerate}

\section{References}
\begin{enumerate}
    \item Stanford University NLP Group. \textit{AgentClinic: a Multimodal Agent Simulation of Clinical Workflows}. GitHub Repository, 2024.
    \item MIT NLP. \textit{MedAgentBench: Benchmark for Evaluative Reasoning in Clinical Agents}. GitHub Repository, 2024.
    \item HL7 International. \textit{HL7 FHIR Release 4 (R4) Specification}. 2019. Available at: \url{https://hl7.org/fhir/R4/}.
    \item U.S. Food and Drug Administration (FDA). \textit{Software as a Medical Device (SaMD): Key Clinical Principles}. 2025.
    \item National Library of Medicine (NLM). \textit{RxNorm, LOINC, and SNOMED CT API Portals}. Available at: \url{https://rxnav.nlm.nih.gov/}.
    \item Apollo Research \& Apollo Health. \textit{Clinical Terminology Mapping and Interoperability Sprints}. 2026.
\end{enumerate}

\end{document}
